\section{The Artin--Schreier theorem}

\begin{lemma} \label{lem:p-power extensions}
\lean{X_pow_sub_C_irreducible_iff_of_prime_pow}
\mathlibok
Let $F$ be a field.
Let $p$ be an odd prime.
Then for $a \in F$, $a$ is not a $p$-th power
if and only if for every positive integer $n$, the polynomial $x^{p^n}-a$ is irreducible over $F$.
(Note that this fails for $p=2$, for example $x^4+4 = (x^2-2x+2)(x^2+2x+2)$.)
\end{lemma}
\begin{proof}
\mathlibok
In Mathlib.
\end{proof}

\begin{lemma} \label{lem:ac-quadratic-adjoins-i}
\lean{quadratic_algebraic_closure_no_i}
\leanok
Let $F$ be a field. Suppose that $F$ contains a square root $i$ of $-1$
and there exists a quadratic extension $K/F$ which is algebraically closed.
Then every element of $F$ is a square.
\end{lemma}
\begin{proof}
Let $a$ be an arbitrary element of $F$.
Since $K$ is algebraically closed, it contains some $b$ for which $b^4 = a$.
The minimal polynomial $f(x)$ of $b$ over $F$ has degree dividing $[K:F] = 2$.
If this degree is 1, then $f(x) = x-b$ and clearly $b^2 \in F$.
If the degree is 2, then $f(0)$ is the product of $b$ and another root of $f(x)$,
which must be one of $ib, -b, -bi$.
Hence $f(0) \in F$ equals $b^2$ times an element of $F$, so $b^2 \in F$.
In both cases, $b^2$ is a square root of $a$ in $F$.
\end{proof}

\begin{lemma} \label{lem:finite-ac-extension-separable}
\lean{finite_algebraic_closure_separable}
\leanok
Let $F$ be a field which admits a finite algebraically closed extension $K$.
Then $K/F$ is a separable extension.
\end{lemma}
\begin{proof}
\uses{lem:ac-quadratic-adjoins-i,lem:p-power extensions}
It suffices to check that $F$ is perfect; this is automatic if $F$ has characteristic $0$, so we
may assume that $F$ is of characteristic $p>0$, and then the perfect condition can be interpreted as the condition that every element of $F$ is a $p$-th power.

Suppose that $p=2$. Since $1 = -1$ in $F$, $F$ contains a square root of $-1$.
We can thus apply Lemma~\ref{lem:ac-quadratic-adjoins-i} to see that every element of $F$ is a perfect square, so $F$ is perfect.

Suppose that $p > 2$. If $F$ is not perfect, we can find $a \in F$ which is not a $p$-th power.
Choose a positive integer $n$ with $p^n > [K:F]$,
then choose $b \in K$ with $b^{p^n} = a$.
The polynomial $x^{p^n} - a$ is irreducible over $F$ by Lemma~\ref{lem:p-power extensions} and has $b$ as a root, so it is the minimal polynomial of $b$ over $F$.
However, the degree of the minimal polynomial is at most $[K:F] < p^n$, a contradiction.
\end{proof}


\begin{proposition} \label{prop:artin schreier extensions}
Let $F$ be a field of prime characteristic $p$.
Let $K/F$ be a Galois extension of degree $p$.
Then $K \cong F[x]/(x^p-x-a)$ for some $a \in F$.
\end{proposition}
\begin{proof}
To be filled in.
\end{proof}


\begin{lemma} \label{lem:cyclic-p-nonzero}
\lean{finite_algebraic_closure_cyclic_prime}
\leanok
Let $F$ be a field which admits an algebraically closed extension $K$ such that $K/F$ is Galois of prime degree $p$. Then $p \neq 0$ in $F$.
\end{lemma}
\begin{proof}
\uses{prop:artin schreier extensions}
If $p=0$ in $F$, then by Proposition~\ref{prop:artin schreier extensions},
$K = F[x]/(x^p-x-a)$ for some $a \in F$. Since $K$ is algebraically closed,
there exists $y \in K$ such that $y^p-y=ax^{p-1}$. Write $y = \sum_{i=0}^{p-1} y_i x^i$ with $y_i \in F$;
then $y_{p-1}^{p} - y_{p-1} = a$ and so $x-y_{p-1} \in \mathbb{F}_p$, a contradiction.
\end{proof}

\begin{lemma} \label{lem:ac-quadratic}
\lean{finite_algebraic_closure_cyclic_quadratic}
\leanok
Let $F$ be a field which admits an algebraically closed extension $K$ such that $K/F$ is Galois of prime degree $p$. Then $p = 2$.
\end{lemma}
\begin{proof}
\uses{lem:cyclic-p-nonzero,lem:p-power extensions}
Suppose by way of contradiction that $p \neq 2$.
By Lemma~\ref{lem:cyclic-p-nonzero}, $F$ cannot have characteristic $p$;
consequently, $K$ contains a primitive $p$-th root of unity $\zeta_p$.
Since $\Q(\zeta_p)/\Q$ is Galois of degree $p-1$, $F(\zeta_p)/F$ is Galois of degree dividing $p-1$;
but this degree must also divide $[K:F] = p$, so it is 1. That is, $\zeta_p \in F$.

By Kummer theory, $K  \cong F[x]/(x^p-a)$ for some $a \in F$. 
Lemma~\ref{lem:p-power extensions} implies that $x^{p^2}-a$ is irreducible over $F$,
so adjoining a root of this polynomial in $K$ gives a subfield of $K$ of degree $p^2$ over $F$.
However, the degree of any subfield must divide $[K:F] = p$, a contradiction.
\end{proof}


\begin{lemma} \label{lem:ac-finite-alg-closure-with-i}
\lean{finite_algebraic_closure_with_i}
\leanok
Let $F$ be a field containing a square root of $-1$ and admitting a finite algebraically closed extension. Then $F$ is algebraically closed.
\end{lemma}
\begin{proof}
\uses{lem:finite-ac-extension-separable,lem:ac-quadratic,lem:finite-ac-extension-separable,lem:ac-quadratic-adjoins-i}
The extension $K/F$ is normal, and by Lemma~\ref{lem:finite-ac-extension-separable} it is separable, so it is Galois; let $G$ be the Galois group. 
If $K \neq F$, then for any prime factor $p$ of $[L:K]$, by Cauchy's theorem 
$G$ contains a nontrivial cyclic subgroup of prime order $p$, whose fixed field is a subfield $E$ of $K$ containing $F$.
Applying Lemma~\ref{lem:ac-quadratic} to $K/E$ yields $p=2$.

By Lemma~\ref{lem:finite-ac-extension-separable}, $F$ is not of characteristic $2$.
We can thus write $K = F[x]/(x^2-a)$ for some $a \in F$ not a perfect square,
but by Lemma~\ref{lem:ac-quadratic-adjoins-i} no such $a$ exists.
\end{proof}

\begin{definition} \label{def:real closed field}
\lean{IsRealClosed}
\mathlibok
A field $F$ is \emph{real} if $-1$ is not a sum of squares in $F$.
A field $F$ is \emph{real closed} is all of the following conditions hold:
\begin{itemize}
\item $F$ is real;
\item for every $x \in F$, one of $x$ or $-x$ is a square in $F$;
\item every odd-degree polynomial over $F$ has a root in $F$.
\end{itemize}
\end{definition}

\begin{lemma} \label{lem:real closed condition}
\lean{RealClosed_from_quadratic}
\leanok
\uses{def:real closed field}
Let $F$ be a field. Suppose that $-1$ is not a square in $F$
and there exists a quadratic extension $K/F$ which is algebraically closed. Then $F$ is real closed.
\end{lemma}
\begin{proof}
Let $i \in K$ be a square root of $-1$.
The sum of any two squares in $F$ is a square: for $a,b \in F$, we can choose $u \in K$ with $u^2 = a+bi$, and then $\Norm_{K/F}(u)^2 = \Norm_{K/F}(a+bi) = a^2+b^2$.
Since $-1$ is not a square in $F$, it follows that $F$ is real.

If $x \in F$ is not a square, then $K=F(x^{1/2})$; by Kummer theory, $-1$ and $x$ belong to the same square class, and so $x^{1/2}/i \in F$ squares to $-x$.

Finally, since $F$ is real, it is of characteristic 0, so $K/F$ is a Galois extension.
For any polynomial $f(x)$ of odd degree over $F$, the Galois group of $K/F$ acts on the roots of $f$ with orbits of order 1 or 2, so there must be a fixed point; hence $f$ has a root in $F$.
\end{proof}

\begin{theorem} \label{thm:artin schreier thm}
\uses{def:real closed field}
\lean{artin_schreier_thm}
\leanok
Let $F$ be a field which admits a finite algebraically closed extension.
Then $F$ is algebraically closed or real closed.
\end{theorem}
\begin{proof}
\uses{lem:ac-finite-alg-closure-with-i,lem:real closed condition}
Let $K$ be an algebraic closure of $F$.
Let $i \in K$ be a square root of $-1$.
By Lemma~\ref{lem:ac-finite-alg-closure-with-i}, $F(i)$ is algebraically closed;
so either $i \in F$ and $F$ is algebraically closed, or $i \notin F$ and 
$F$ is real closed by Lemma~\ref{lem:real closed condition}.
\end{proof}

